% main_calculus_lecture2_metropolis.tex
\documentclass[aspectratio=169]{beamer}

% --- Language & encoding ---
\usepackage[utf8]{inputenc}
\usepackage[T1]{fontenc}
\usepackage[english,french]{babel}
\usepackage{lmodern}
\usepackage{microtype}

% --- Math & tables ---
\usepackage{amsmath,amssymb}
\usepackage{booktabs}
\usepackage{graphicx}
\usepackage{changepage}
\usepackage{amssymb}
\newcommand{\R}{\mathbb{R}}
% --- Metropolis Theme ---
\usetheme[numbering=fraction,progressbar=frametitle]{metropolis}
\setbeamertemplate{frame numbering}[fraction]

% --- Colors (keep your palette) ---
\definecolor{BootcampBlueDark}{HTML}{1A3C68}
\definecolor{TextBlack}{HTML}{000000}
\definecolor{AccentGold}{HTML}{DAA520}
\definecolor{Crimson}{HTML}{990000}
\definecolor{MatrixGreen}{HTML}{228B22}
\definecolor{MatrixRed}{HTML}{B22222}

\setbeamercolor{block title example}{bg=BootcampBlueDark!90!black,fg=white}
\setbeamercolor{block body example}{bg=BootcampBlueDark!5!white,fg=black}

% === Thick frame title band ===
\setbeamercolor{frametitle}{bg=BootcampBlueDark,fg=white}
\setbeamerfont{frametitle}{series=\bfseries,size=\large}
\setbeamertemplate{frametitle}{
  \nointerlineskip
  \begin{beamercolorbox}[wd=\paperwidth,ht=3.2ex,dp=1.4ex,leftskip=1em,rightskip=1em]{frametitle}
    \usebeamerfont{frametitle}\insertframetitle
  \end{beamercolorbox}
}

% === Custom Definition Block ===
\newenvironment{definitionblock}[1]{%
  \par\vspace{0.4em}%
  \begin{beamercolorbox}[
      wd=\textwidth,
      sep=0.4em,
      leftskip=0.3em,
      rightskip=0.6em,
      rounded=false,
      shadow=false
    ]{block title example}%
    \raisebox{0.15em}{\usebeamerfont{block title example}#1}%
  \end{beamercolorbox}%
  \vspace{-0.4em}%
  \begin{beamercolorbox}[
      wd=\textwidth,
      sep=1em,
      leftskip=0.6em,
      rightskip=0.6em,
      rounded=false,
      shadow=false
    ]{block body example}%
    \usebeamerfont{block body example}%
}{%
  \end{beamercolorbox}%
  \par\vspace{0.6em}%
}

% === Global text formatting ===
\setbeamercolor{normal text}{fg=TextBlack,bg=white}
\setbeamercolor{title}{fg=TextBlack}
\setbeamercolor{structure}{fg=BootcampBlueDark}
\setbeamercolor{progress bar}{fg=BootcampBlueDark,bg=gray!30}

% --- Course Info ---
\title{Fourrier Analysis}
\institute{NYU Tandon School, Finance and Risk Engineering}

\metroset{
  sectionpage=progressbar,
  subsectionpage=progressbar
}

% --- Section & subsection transitions (plain) ---
\AtBeginSection{\frame[plain]{\sectionpage}}
\AtBeginSubsection{\frame[plain]{\subsectionpage}}

\begin{document}

\maketitle

\begin{frame}{Outline}
  \tableofcontents[hideallsubsections]
\end{frame}

% ==================================================
% SECTION 1 — Differential Calculus (1D → nD)
% ==================================================
\section{Introduction}

\begin{frame}{Outline}
  \tableofcontents[currentsection,hideothersubsections]
\end{frame}

% ======================================================
\subsection{Fourrier Analysis : Motivation}
% ======================================================

\begin{frame}{Harmonic Analysis}

Fourier analysis is the study of how a signal can be decomposed into elementary oscillations.
\pause
\begin{itemize}
    \item Many phenomena are \textbf{periodic} or have hidden periodic structure:
    \begin{itemize}
        \item sound, vibrations, electromagnetic waves
        \item temperature cycles, seasonal components
        \item financial time series (volatility clustering, periodicities)
    \end{itemize}
\pause
    \item Harmonic analysis represents a function as a \textbf{combination of sines and cosines}
    \[
        f(t) \approx \sum_{n=-N}^{N} c_n e^{i n \omega t}
    \]
\pause
    \item Key idea:
    \begin{block}{Decomposition principle}
      Complex signals are built from simple oscillations.
    \end{block}
\end{itemize}

\end{frame}

\begin{frame}{Signal Treatment and Machine Learning}

\begin{itemize}
    \item Many data types can be viewed as \textbf{signals}:
    \begin{itemize}
        \item time series
        \item images
        \item text (word frequency)
    \end{itemize}
\pause
    \item Fourier analysis allows:
    \begin{itemize}
        \item noise reduction (filtering high-frequency noise)
        \item feature extraction
        \item spectral analysis
        \item denoising and compression
    \end{itemize}
\pause
    \item In finance:
    \begin{itemize}
        \item analysis of \textbf{cycles} in returns
        \item volatility models in frequency domain
        \item option pricing via characteristic functions
    \end{itemize}
\pause
    \item In machine learning:
    \begin{itemize}
        \item convolutional neural networks use frequency filters
        \item Fourier features for kernel methods
    \end{itemize}
\end{itemize}

\end{frame}


% ======================================================
\subsection{The two main paradgigm}
% ======================================================

\begin{frame}{Fourier Series}

\begin{block}{Goal}
Represent a \textbf{periodic} function as a sum of sines and cosines.
\end{block}
\pause
\[
f(t) = \sum_{n=-\infty}^{\infty} c_n e^{i n \omega_0 t}
\]

with coefficients
\[
c_n = \frac{1}{T} \int_{0}^{T} f(t)e^{-in\omega_0 t}\, dt
\]
\pause
\begin{itemize}
    \item Works only for \textbf{periodic} signals
    \item Frequency spectrum is \textbf{discrete}
    \item Natural tool for:
    \begin{itemize}
        \item PDEs on bounded domains
        \item vibration modes
        \item periodic time series
    \end{itemize}
\end{itemize}

\end{frame}

\begin{frame}{Fourier Transform}

\begin{block}{Goal}
Represent a \textbf{non-periodic} function using a continuous spectrum of frequencies.
\end{block}
\pause
\[
\hat{f}(\omega)=\int_{-\infty}^{\infty} f(t)\, e^{-i\omega t}\,dt
\]

\[
f(t)=\frac{1}{2\pi}\int_{-\infty}^{\infty} \hat{f}(\omega)\, e^{i\omega t}\, d\omega
\]
\pause
\begin{itemize}
    \item Works for non-periodic signals
    \item Frequency spectrum is \textbf{continuous}
    \item Used for:
    \begin{itemize}
        \item signal processing
        \item probability characteristic functions
        \item option pricing (Fourier inversion)
        \item PDE solution (heat equation)
    \end{itemize}
\end{itemize}
\pause
\end{frame}

\section{Fourrier Series}

\subsection{One -Dimensional}

\subsubsection{Definition}

\begin{frame}{Definition of Fourier Series}

Let \(f\) be a periodic function with period \(T\).

\[
f(t+T)=f(t)
\]
\pause
The Fourier series of \(f\) is

\[
f(t) \sim \frac{a_0}{2}
+ \sum_{n=1}^{\infty}
\left(
a_n \cos\left(\frac{2\pi n t}{T}\right)
+
b_n \sin\left(\frac{2\pi n t}{T}\right)
\right)
\]
\pause
\begin{itemize}
    \item \(a_n\): cosine coefficients (even part)
    \item \(b_n\): sine coefficients (odd part)
    \item Frequency \(n/T\)
\end{itemize}

\end{frame}

\begin{frame}{Fourier Coefficients}

The coefficients are computed as projections on the orthogonal basis.

\[
a_n = \frac{2}{T}\int_{0}^{T}
f(t)\cos\left(\frac{2\pi n t}{T}\right)\,dt
\]
\pause
\[
b_n = \frac{2}{T}\int_{0}^{T}
f(t)\sin\left(\frac{2\pi n t}{T}\right)\,dt
\]
\pause
\[
a_0 = \frac{2}{T}\int_{0}^{T} f(t)\,dt
\]

\begin{block}{Interpretation}
Each coefficient measures ``how much'' of each oscillation is present in the signal.
\end{block}

\end{frame}

\begin{frame}{Complex Fourier Series}

Using $
e^{i\theta}=\cos\theta + i\sin\theta
$, we can write
\[
f(t)=\sum_{n=-\infty}^{\infty} c_n e^{i 2\pi n t/T}
\]
with $
c_n=\frac{1}{T}\int_{0}^{T}
f(t)e^{-i 2\pi n t/T}\, dt
.$
\pause
\begin{itemize}
    \item spectrum is \textbf{discrete}
    \item natural in probability and finance
    \item directly linked to \textbf{characteristic functions}
\end{itemize}

\end{frame}

\subsubsection{Properties}

\begin{frame}{Orthogonality relations}

On the interval $[-\pi,\pi]$:

\[
\int_{-\pi}^{\pi} \cos(nx)\cos(mx)\,dx =
\begin{cases}
0 & n\neq m \\
\pi & n=m\neq 0
\end{cases}
\]
\pause
\[
\int_{-\pi}^{\pi} \sin(nx)\sin(mx)\,dx =
\begin{cases}
0 & n\neq m \\
\pi & n=m
\end{cases}
\]
\pause
\[
\int_{-\pi}^{\pi} \sin(nx)\cos(mx)\,dx = 0
\]
\pause
\begin{block}{Interpretation}
Sines and cosines form an orthogonal basis of $L^2[-\pi,\pi]$.
\end{block}

\end{frame}

\begin{frame}{Linearity}

If
\[
f(x) \sim \sum (a_n \cos nx + b_n \sin nx)
\]
and
\[
g(x) \sim \sum (\alpha_n \cos nx + \beta_n \sin nx)
\]

then for any constants $\lambda,\mu$:

\[
\lambda f(x)+\mu g(x)
\]
has Fourier coefficients

\[
a_n^{(\lambda f+\mu g)}=\lambda a_n+\mu \alpha_n
\qquad
b_n^{(\lambda f+\mu g)}=\mu b_n+\mu \beta_n
\]

\begin{block}{Important}
Fourier transform is a linear operator.
\end{block}

\end{frame}

\begin{frame}{Even and odd functions}

\textbf{If $f$ is even}: $
f(-x)=f(x)
$
then:
\[
b_n = 0
\]
Only cosines appear.
\pause
\textbf{If $f$ is odd}:
$
f(-x)=-f(x)
$
then:
\[
a_n = 0,\quad a_0=0
\]
Only sines appear.
\begin{block}{Practical consequence}
Huge simplification of calculations.
\end{block}

\end{frame}

\begin{frame}{Effect of translation}

If

\[
g(x)=f(x-x_0)
\]

then Fourier coefficients satisfy

\[
c_n(g)=c_n(f)e^{-inx_0}
\]
\pause
\begin{itemize}
    \item translation in time
    \item corresponds to phase shift in frequency domain
\end{itemize}

\begin{block}{Signal processing interpretation}
Shifting a signal does not change the amplitude spectrum,
only the phase.
\end{block}

\end{frame}

\begin{frame}{Scaling (time dilation) for Fourier series}

Let $f$ be a $2\pi$–periodic function with Fourier series

\[
f(x)=\sum_{n=-\infty}^{\infty} c_n e^{inx}.
\]

Define the scaled function
\[
g(x)=f(ax), \qquad a\neq 0.
\]
\pause
\begin{block}{Period change}
If $f$ has period $2\pi$, then $g$ has period
\[
\text{new period }=\frac{2\pi}{|a|}.
\]
\end{block}

\begin{block}{Effect on Fourier series}
\[
g(x)=f(ax)
=
\sum_{n=-\infty}^{\infty}
c_n\,e^{ina x}.
\]

So the frequencies are multiplied by $a$.
\end{block}

\begin{block}{Intuition}
Compressing the signal in time $\Rightarrow$ increases frequencies.  
Stretching the signal in time $\Rightarrow$ decreases frequencies.
\end{block}

\end{frame}

\begin{frame}{Differentiation and integration}

If the Fourier series converges sufficiently well:

\[
f(x)=\sum a_n\cos(nx)+b_n\sin(nx)
\]

then

\[
f'(x)=\sum -n a_n \sin(nx)+ n b_n \cos(nx)
\]
\pause
Implications:

\begin{itemize}
    \item differentiation multiplies coefficients by \(n\)
    \item integration divides coefficients by \(n\)
\end{itemize}

\begin{block}{Risk engineering remark}
High frequencies are amplified by differentiation:
→ numerical instability, noise amplification.
\end{block}

\end{frame}

\begin{frame}{Parseval Identity}

Energy in time domain = energy in frequency domain

\[
\frac{1}{\pi}\int_{-\pi}^{\pi} |f(x)|^2 \, dx
=
\frac{a_0^2}{2}+\sum_{n=1}^{\infty}(a_n^2+b_n^2)
\]
\pause
\begin{block}{Meaning}
The Fourier coefficients describe the distribution of energy among frequencies.
\end{block}

Applications:

\begin{itemize}
    \item variance of time series
    \item spectral density
    \item risk decomposition by frequency
\end{itemize}

\end{frame}

\subsubsection{Examples}

\begin{frame}{Example 1: Fourier series of $f(x)=x$ on $[-\pi,\pi]$}

We consider the $2\pi$-periodic function

\[
f(x) = x,\quad x \in [-\pi,\pi].
\]
\pause
\begin{itemize}
    \item $f$ is an \textbf{odd} function: $f(-x)=-f(x)$
    \item $\Rightarrow$ Fourier series contains only \textbf{sine} terms
\end{itemize}
\pause
We want to compute its Fourier series:

\[
f(x) \sim \sum_{n=1}^{\infty} b_n \sin(nx).
\]

\end{frame}

\begin{frame}{Example 1: computation of the coefficients}

Since $f$ is odd:
$
a_0 = 0, \quad a_n = 0.
$
\pause
We only need $b_n$:

\[
b_n = \frac{1}{\pi}\int_{-\pi}^{\pi} x \sin(nx)\,dx.
\]

$f$ is odd, $\sin(nx)$ is odd, so $x\sin(nx)$ is \textbf{even}:
$
b_n = \frac{2}{\pi}\int_{0}^{\pi} x \sin(nx)\,dx.
$
\pause
Integrate by parts:
\[
\begin{aligned}
u &= x, & dv &= \sin(nx)\,dx \\
du &= dx, & v &= -\frac{1}{n}\cos(nx)
\end{aligned}
\]
\pause
\[
\int_{0}^{\pi} x \sin(nx)\,dx
=
\left[-\frac{x}{n}\cos(nx)\right]_{0}^{\pi}
+ \frac{1}{n}\int_{0}^{\pi}\cos(nx)\,dx.
\]

\end{frame}

\begin{frame}{Example 1: final expression}

We have $
\int_{0}^{\pi}\cos(nx)\,dx
= \left[\frac{1}{n}\sin(nx)\right]_{0}^{\pi}=0.
$
\pause
So

$$
\int_{0}^{\pi} x \sin(nx)\,dx
= -\frac{\pi}{n}\cos(n\pi) - 0
= -\frac{\pi}{n}(-1)^n.
$$
\pause
Thus

\[
b_n = \frac{2}{\pi}\left(-\frac{\pi}{n}(-1)^n\right)
= \frac{2}{\pi}\cdot\left(-\frac{\pi}{n}(-1)^n\right)
= \frac{2}{n}(-1)^{n+1}.
\]
\pause
\begin{block}{Fourier series of $x$ on $[-\pi,\pi]$}
\[
x \sim 2\sum_{n=1}^{\infty} \frac{(-1)^{n+1}}{n}\sin(nx).
\]
\end{block}

\end{frame}

\begin{frame}{Example 2: Fourier series of $f(x)=|x|$ on $[-\pi,\pi]$}

We consider

\[
f(x)=|x|,\quad x\in[-\pi,\pi],
\]

extended as a $2\pi$-periodic function.
\pause
\begin{itemize}
    \item $f$ is \textbf{even}: $f(-x)=f(x)$
    \item $\Rightarrow$ Fourier series contains only \textbf{cosine} terms
\end{itemize}

So

\[
|x| \sim \frac{a_0}{2} + \sum_{n=1}^{\infty} a_n \cos(nx),
\quad b_n = 0.
\]

\end{frame}

\begin{frame}{Example 2: computation of $a_0$}

\[
a_0 = \frac{1}{\pi}\int_{-\pi}^{\pi} |x|\,dx.
\]
\pause
Since $|x|$ is even:

\[
a_0 = \frac{2}{\pi}\int_{0}^{\pi} x\,dx
    = \frac{2}{\pi}\left[\frac{x^2}{2}\right]_{0}^{\pi}
    = \frac{2}{\pi}\cdot \frac{\pi^2}{2}
    = \pi.
\]
\pause
So

\[
\frac{a_0}{2} = \frac{\pi}{2}.
\]

\end{frame}

\begin{frame}{Example 2: computation of $a_n$}

For $n\ge 1$: $
a_n = \frac{1}{\pi}\int_{-\pi}^{\pi} |x|\cos(nx)\,dx.
$ Again using evenness:

\[
a_n = \frac{2}{\pi}\int_{0}^{\pi} x\cos(nx)\,dx.
\]
\pause
Integrate by parts:

\[
\begin{aligned}
u &= x, & dv &= \cos(nx)\,dx \\
du &= dx, & v &= \frac{1}{n}\sin(nx)
\end{aligned}
\]
\pause
\[
\int_{0}^{\pi} x\cos(nx)\,dx
=
\left[\frac{x}{n}\sin(nx)\right]_{0}^{\pi}
- \frac{1}{n}\int_{0}^{\pi}\sin(nx)\,dx.
\]
\pause
\[
\int_{0}^{\pi}\sin(nx)\,dx
=
\left[-\frac{1}{n}\cos(nx)\right]_{0}^{\pi}
=
-\frac{1}{n}(\cos(n\pi)-1).
\]

\end{frame}

\begin{frame}{Example 2: final expression}

First note that $\sin(n\pi)=0$, so the boundary term is $0$: $
\left[\frac{x}{n}\sin(nx)\right]_{0}^{\pi} = 0
$. Thus
\[
\int_{0}^{\pi} x\cos(nx)\,dx
= -\frac{1}{n}\left(-\frac{1}{n}(\cos(n\pi)-1)\right)
= \frac{1}{n^2}(\cos(n\pi)-1).
\]
\pause
So

\[
a_n = \frac{2}{\pi}\cdot \frac{1}{n^2}(\cos(n\pi)-1)
= \frac{2}{\pi n^2}\big((-1)^n-1\big).
\]

\begin{itemize}
    \item If $n$ is even: $(-1)^n=1 \Rightarrow a_n=0$.
    \item If $n$ is odd: $(-1)^n=-1 \Rightarrow a_n=\dfrac{2}{\pi n^2}(-2)=-\dfrac{4}{\pi n^2}$.
\end{itemize}

\end{frame}

\begin{frame}{Example 2: final Fourier series of $|x|$}

Only odd indices survive. Let $n=2k+1$:

\[
a_{2k+1} = -\frac{4}{\pi (2k+1)^2},\quad k=0,1,2,\dots
\]
\pause
Therefore

\[
|x| \sim \frac{\pi}{2}
-\frac{4}{\pi}\sum_{k=0}^{\infty}
\frac{\cos\big((2k+1)x\big)}{(2k+1)^2},
\quad x\in[-\pi,\pi].
\]
\pause
\begin{block}{Remark}
This series is useful to derive interesting identities such as
\[
1 + \frac{1}{3^2} + \frac{1}{5^2} + \cdots = \frac{\pi^2}{8}.
\]
\end{block}

\end{frame}

\begin{frame}{How to compute a Fourier series: general method}

Goal: find the Fourier expansion of a $2\pi$-periodic function $f$.
\pause
\begin{block}{Step 1 — Identify if there is a symmetry}
\begin{itemize}
    \item If $f$ is \textbf{even} ($f(-x)=f(x)$): $
    b_n = 0 \quad \Rightarrow \text{cosine series only}
    $
    \item If $f$ is \textbf{odd} ($f(-x)=-f(x)$):
    $
    a_n = 0,\; a_0=0 \quad \Rightarrow \text{sine series only}
    $
\end{itemize}
\end{block}
\pause
\vspace{-5mm}
\begin{block}{Step 2 — Choose the correct interval}
Most common: $
[-\pi,\pi] \quad \text{or} \quad [0,2\pi]
$, adapt formulas to period $T$ if necessary.
\end{block}
\pause
\begin{block}{Step 3 — Compute coefficients}
$
a_0=\frac{1}{\pi}\int_{-\pi}^{\pi} f(x)\,dx
$, $
a_n=\frac{1}{\pi}\int_{-\pi}^{\pi} f(x)\cos(nx)\,dx
$ and, $
b_n=\frac{1}{\pi}\int_{-\pi}^{\pi} f(x)\sin(nx)\,dx
$
Use symmetry to simplify integrals.
\end{block}
\pause
\begin{block}{Step 4 — Write the series}
\[
f(x)\sim \frac{a_0}{2}+\sum_{n=1}^{\infty}
\left(a_n\cos nx + b_n\sin nx\right)
\]
State where it converges (continuous points / jump midpoints).
\end{block}

\end{frame}

\subsubsection{Remark}


\begin{frame}{General period $T$ (not only $2\pi$)}

The Fourier method works exactly the same for any period $T>0$.

\begin{block}{Periodicity}
\vspace{-5mm}
\[
f(t+T)=f(t)
\vspace{-5mm}
\]
\end{block}
\pause
\begin{block}{Fundamental frequency}
\vspace{-5mm}
\[
\omega_0=\frac{2\pi}{T}
\vspace{-5mm}
\]
\end{block}
\pause
\begin{block}{Fourier series with general period $T$}
\vspace{-2mm}
\[
f(t)\sim
\frac{a_0}{2}
+\sum_{n=1}^{\infty}
\left(
a_n\cos(n\omega_0 t)+
b_n\sin(n\omega_0 t)
\right)
\vspace{-9mm}
\]
\end{block}
\pause
\begin{block}{Coefficients}
\[
a_0=\frac{2}{T}\int_{0}^{T} f(t)\,dt
\,\, , \,\, 
a_n=\frac{2}{T}\int_{0}^{T} f(t)\cos(n\omega_0 t)\,dt
\,\, , \,\,
b_n=\frac{2}{T}\int_{0}^{T} f(t)\sin(n\omega_0 t)\,dt
\vspace{-7mm}
\]
\end{block}
\pause
\begin{block}{Key message}
Same method as for $2\pi$: $
\textbf{only the frequency } \omega_0=\frac{2\pi}{T} \text{ changes.}$
\end{block}
\end{frame}


\begin{frame}{Convergence of Fourier Series (intuition)}

\begin{itemize}
    \item If \(f\) is continuous and piecewise smooth:
    \[
    \text{Fourier series converges to } f
    \]
    \item If \(f\) has jumps:
    \[
    \text{converges to the midpoint of the jump}
    \]
    \item Convergence is in:
    \begin{itemize}
        \item pointwise
        \item \(L^2\) sense (energy)
    \end{itemize}
\end{itemize}
\pause
\begin{block}{Dirichlet Conditions}
Finite number of maxima, minima, and discontinuities on each period is sufficient.
\end{block}

\end{frame}

\begin{frame}{Applications in Financial Risk Engineering}

\begin{itemize}
    \item Spectral analysis of time series
    \item Seasonal adjustment of returns
    \item Noise filtering of price signals
    \item Option pricing:
    \begin{itemize}
        \item Fourier inversion of characteristic functions
        \item Carr-Madan method
        \item Heston model solutions
    \end{itemize}
    \item Estimation of volatility via frequency domain
\end{itemize}

\end{frame}


\subsection{d-dimensional}

\begin{frame}{Why Fourier series in several dimensions?}

We often study functions of several variables:
\[
f(x_1,\dots,x_d)=f(x)
\]

Examples:
\begin{itemize}
    \item temperature depending on space $(x,y,z)$
    \item images (2D signals)
    \item solutions of PDEs in several dimensions
    \item financial models with several assets
\end{itemize}
\pause
Goal:
\begin{block}{Multidimensional harmonic analysis}
Decompose $f$ as a combination of oscillations in all directions.
\end{block}

\end{frame}

\begin{frame}{Periodicity in $\mathbb{R}^d$}

We work on a $d$-dimensional box:
\[
Q = [0,L_1]\times \cdots \times [0,L_d].
\]
\pause
A function $f:\mathbb{R}^d\to\mathbb{C}$ is periodic if
\[
f(x+L_j e_j)=f(x)\qquad \text{for all } j=1,\dots,d,
\]
where $e_j$ is the $j$-th coordinate vector.
\pause
\begin{block}{Interpretation}
$f$ repeats itself along each coordinate direction.
\end{block}

\end{frame}

\begin{frame}{Multidimensional Fourier basis}

Define the fundamental frequencies
\[
\omega_j = \frac{2\pi}{L_j}, \qquad j=1,\dots,d.
\]
\pause
For $k=(k_1,\dots,k_d)\in\mathbb{Z}^d$ define
\[
e_k(x)=\exp\left(i\sum_{j=1}^d k_j \omega_j x_j\right)
= e^{\,i\,k\cdot \omega \cdot x}.
\]
\pause
\begin{block}{Key fact}
The family $\{e_k\}_{k\in\mathbb{Z}^d}$
is an orthogonal basis of $L^2(Q)$.
\end{block}

This generalizes $\{e^{inx}\}$ in dimension 1.
\end{frame}

\begin{frame}{Fourier series in dimension $d$}

Any sufficiently regular periodic function $f$ can be written as
\[
f(x)=\sum_{k\in\mathbb{Z}^d} c_k \, e^{\,i\sum_{j=1}^d k_j \omega_j x_j}.
\]
\pause
The Fourier coefficients are
\[
c_k = \frac{1}{|Q|}
\int_{Q} f(x)\,
e^{\,-i\sum_{j=1}^d k_j \omega_j x_j}\,dx,
\]

with volume
\[
|Q| = L_1\cdots L_d.
\]
\pause
\begin{block}{Same structure as 1D}
\[
\text{Projection coefficient = average scalar product}.
\]
\end{block}

\end{frame}

\begin{frame}{Sine–cosine representation}

In two dimensions ($d=2$), for example:

\[
f(x,y)=
\sum_{m,n\in\mathbb{Z}}
c_{mn}\,e^{i(m\omega_1 x + n\omega_2 y)}.
\]

Equivalently:

\[
f(x,y)=\sum_{m,n\ge0}
\Big(
a_{mn}\cos(m\omega_1 x)\cos(n\omega_2 y)
+\cdots
\Big)
\]
\pause
\begin{itemize}
    \item useful for PDEs on rectangles
    \item used in image compression (2D cosine transform)
\end{itemize}

\end{frame}

\begin{frame}{Parseval identity in dimension $d$}

Orthogonality:
\[
\int_Q e_k(x)\,\overline{e_\ell(x)}\,dx
=
\begin{cases}
|Q| & \text{if } k=\ell,\\
0 & \text{otherwise.}
\end{cases}
\]
\pause
Energy equality (Parseval):
\[
\frac{1}{|Q|}\int_Q |f(x)|^2\,dx
=
\sum_{k\in\mathbb{Z}^d} |c_k|^2.
\]

\begin{block}{Meaning}
Total energy/variance = sum of squared Fourier modes.
\end{block}

\end{frame}

\begin{frame}{Financial interpretation (Risk Engineering)}

Multidimensional Fourier series apply to:

\begin{itemize}
    \item joint density of several assets
    \item pricing with several risk factors
    \item dependence in baskets or indices
    \item PDEs in multiple state variables
    \item volatility surface approximation
\end{itemize}

\begin{block}{Example}
Basket option payoff depending on $(S_1,\dots,S_d)$.
\end{block}

Fourier analysis decomposes complexity into frequencies.
\end{frame}

\begin{frame}{Example in 2D}

Let $f(x,y)=x+y$ on $[-\pi,\pi]^2$,
extended periodically.

\begin{itemize}
    \item odd in each variable
    \item Fourier series contains only sine terms
\end{itemize}

We obtain

\[
f(x,y)=
4\sum_{m=1}^{\infty}\sum_{n=1}^{\infty}
\frac{(-1)^{m+n}}{mn}
\sin(mx)\sin(ny).
\]

\begin{block}{Remark}
This is just the tensor product of 1D expansions.
\end{block}

\end{frame}


\section{Fourier Transform}

\begin{frame}{Why Fourier Transform?}

Fourier series apply to periodic functions.
\pause
Many real signals are \textbf{not periodic}:
\begin{itemize}
    \item financial returns
    \item probability densities
    \item impulse responses
    \item localized shocks (events, jumps)
\end{itemize}
\pause
\begin{block}{Idea}
Represent $f(t)$ as a superposition of infinitely many frequencies,
now with a \textbf{continuous} spectrum.
\end{block}

This leads to the Fourier transform.
\end{frame}

\subsection{One-dimensional}

\begin{frame}{Fourier Transform in 1D}

Let $f:\mathbb{R}\to\mathbb{C}$.

\begin{block}{Fourier transform}
\[
\mathcal{F}\{f\}(\xi)
=
\hat{f}(\xi)
=
\int_{-\infty}^{\infty}
f(x)e^{-2\pi i x \xi}\,dx
\]
\end{block}
\pause
\begin{itemize}
    \item spectrum is continuous in $\xi$
    \item same idea as Fourier series but on the whole line
\end{itemize}

\end{frame}

\begin{frame}{Interpretation}

\begin{itemize}
    \item $\hat{f}(\xi)$ measures the presence of frequency $\xi$
    \item $f$ localized in time → $\hat f$ spread in frequency
    \item $f$ oscillatory → $\hat f$ concentrated in frequency
\end{itemize}

\begin{block}{Heisenberg-type tradeoff}
Cannot be localized sharply in both time and frequency.
\end{block}

Applications:
\begin{itemize}
    \item signal processing
    \item machine learning kernels
    \item probability characteristic functions
\end{itemize}

\end{frame}

\begin{frame}{Main properties in 1D}

\begin{itemize}
    \item Linearity:
    \[
    \mathcal{F}(af+bg)=a\hat f+b\hat g
    \]

    \item Translation:
    \[
    f(x-x_0)\quad\Longrightarrow\quad 
    e^{-2\pi i\xi x_0}\hat f(\xi)
    \]

    \item Modulation:
    \[
    e^{2\pi i x \xi_0}f(x)
    \quad\Longrightarrow\quad
    \hat f(\xi-\xi_0)
    \]

    \item Dilation:
    \[
    f(ax)\quad\Longrightarrow\quad
    \frac{1}{|a|}\hat f\!\left(\frac{\xi}{a}\right)
    \]

    \item Differentiation:
    \[
    f'(x)\quad\Longrightarrow\quad
    (2\pi i\xi)\hat f(\xi)
    \]
\end{itemize}

\end{frame}

\begin{frame}{Parseval/Plancherel identity}

\[
\int_{-\infty}^{\infty} |f(x)|^2\,dx
=
\int_{-\infty}^{\infty} |\hat f(\xi)|^2\,d\xi
\]
\pause
\begin{block}{Meaning}
Energy / variance is preserved by the Fourier transform.
\end{block}

\begin{itemize}
    \item orthonormality in $L^2(\mathbb{R})$
    \item frequency-domain risk decomposition
\end{itemize}

\end{frame}

\begin{frame}{Fourier inversion formula}

Let $f:\mathbb{R}\to\mathbb{C}$ such that $f$ and $\hat f$ are integrable
(or more generally $f\in L^2(\mathbb{R})$).

\begin{block}{Fourier transform}
\[
\hat f(\xi)=\int_{-\infty}^{\infty}
f(x)e^{-2\pi i x \xi}\,dx
\]
\end{block}
\pause
\begin{block}{Inversion formula}
Then $f$ can be recovered from $\hat f$:
\[
f(x)=\int_{-\infty}^{\infty}
\hat f(\xi)\,e^{2\pi i x \xi}\,d\xi
\]
\end{block}
\pause
\vspace{-6mm}
\begin{block}{Key message}
No information is lost: the Fourier transform is reversible.
\end{block}

\begin{itemize}
    \item same formula in $\mathbb{R}^d$ with $x\cdot \xi$
    \item inversion holds in the $L^2$ sense (Plancherel theorem)
    \item when $f$ has jumps, inversion gives midpoint values
\end{itemize}

\end{frame}

\begin{frame}{Financial interpretation in 1D}

Let $X$ be a random variable.

\[
\phi_X(\xi)=\mathbb{E}\left[e^{i\xi X}\right]
\]
\pause
\begin{itemize}
    \item $\phi_X$ is the \textbf{characteristic function}
    \item it is essentially a Fourier transform of the density
\end{itemize}

Applications:
\begin{itemize}
    \item Carr–Madan formula
    \item Heston model
    \item Variance gamma, Lévy models
\end{itemize}

\end{frame}

\begin{frame}{Example 1: Fourier transform of an indicator}

Consider
\[
f(x)=\mathbf{1}_{[-1,1]}(x)=
\begin{cases}
1 & \text{if } x\in[-1,1],\\
0 & \text{otherwise.}
\end{cases}
\]
\pause
We compute its Fourier transform:
\[
\hat f(\xi)=\int_{-\infty}^{\infty}
f(x)e^{-2\pi i x\xi}\,dx.
\]
\pause
Since $f$ is zero outside $[-1,1]$:
\[
\hat f(\xi)=\int_{-1}^{1}
e^{-2\pi i x\xi}\,dx.
\]

\end{frame}

\begin{frame}{Example 1: explicit computation}

We have
\[
\hat f(\xi)=\int_{-1}^{1}
e^{-2\pi i x\xi}\,dx.
\]
\pause
For $\xi\neq 0$:
\[
\int e^{-2\pi i x\xi}\,dx
=
\frac{1}{-2\pi i\xi}e^{-2\pi i x\xi}.
\]
\pause
So
\[
\hat f(\xi)
=
\left[
\frac{1}{-2\pi i\xi}e^{-2\pi i x\xi}
\right]_{x=-1}^{x=1}
=
\frac{1}{-2\pi i\xi}
\left(
e^{-2\pi i \xi}-e^{2\pi i \xi}
\right).
\]

\end{frame}

\begin{frame}{Example 1: final form (sinc function)}

Recall
\[
e^{-i\theta}-e^{i\theta} = -2i\sin\theta.
\]

With $\theta=2\pi\xi$:
\[
e^{-2\pi i \xi}-e^{2\pi i \xi}
= -2i\sin(2\pi\xi).
\]
\pause
Thus
\[
\hat f(\xi)
=
\frac{1}{-2\pi i\xi}\big(-2i\sin(2\pi\xi)\big)
=
\frac{\sin(2\pi\xi)}{\pi\xi},
\quad \xi\neq 0.
\]
\pause
At $\xi=0$, we have
\[
\hat f(0)=\int_{-1}^{1}1\,dx=2,
\]
and
\[
\lim_{\xi\to 0}\frac{\sin(2\pi\xi)}{\pi\xi}=2.
\]

\begin{block}{Result}
\[
\boxed{\hat f(\xi)=\frac{\sin(2\pi\xi)}{\pi\xi}}
\quad \text{(with the continuous extension } \hat f(0)=2\text{).}
\]
\end{block}

\end{frame}

\begin{frame}{Example 2: Fourier transform of a Laplace-type function}

Let $a>0$ and
\[
f(x)=e^{-a|x|},\quad x\in\mathbb{R}.
\]

We want
\[
\hat f(\xi)=\int_{-\infty}^{\infty}
e^{-a|x|}e^{-2\pi i x\xi}\,dx.
\]
\pause
Note:
\begin{itemize}
    \item $f$ is even: $f(-x)=f(x)$
    \item $\Rightarrow \hat f(\xi)$ will be real and even.
\end{itemize}

\end{frame}

\begin{frame}{Example 2: splitting the integral}

We split according to $|x|$:

\[
\hat f(\xi)
=
\int_{-\infty}^{0} e^{-a(-x)}e^{-2\pi i x\xi}\,dx
+
\int_{0}^{\infty} e^{-ax}e^{-2\pi i x\xi}\,dx.
\]
\pause
So
\[
\hat f(\xi)
=
\int_{-\infty}^{0} e^{ax}e^{-2\pi i x\xi}\,dx
+
\int_{0}^{\infty} e^{-ax}e^{-2\pi i x\xi}\,dx.
\]
\pause
Change variable in the first integral $x\mapsto -x$:

\[
\int_{-\infty}^{0} e^{ax}e^{-2\pi i x\xi}\,dx
=
\int_{0}^{\infty} e^{-ax}e^{2\pi i x\xi}\,dx.
\]
\pause
Thus
\[
\hat f(\xi)
=
\int_{0}^{\infty} e^{-ax}e^{2\pi i x\xi}\,dx
+
\int_{0}^{\infty} e^{-ax}e^{-2\pi i x\xi}\,dx.
\]

\end{frame}

\begin{frame}{Example 2: using cosine representation}

Factor $e^{-ax}$:

\[
\hat f(\xi)
=
\int_{0}^{\infty} e^{-ax}
\left(e^{2\pi i x\xi}+e^{-2\pi i x\xi}\right)\,dx.
\]
\pause
But
\[
e^{2\pi i x\xi}+e^{-2\pi i x\xi}=2\cos(2\pi x\xi).
\]
\pause
So
\[
\hat f(\xi)
=2\int_{0}^{\infty} e^{-ax}\cos(2\pi x\xi)\,dx.
\]
\pause
We now use the standard integral (for $a>0$):
\[
\int_{0}^{\infty} e^{-ax}\cos(bx)\,dx
=
\frac{a}{a^2+b^2}.
\]

Here $b=2\pi\xi$.
\end{frame}

\begin{frame}{Example 2: final expression}

Using
\[
\int_{0}^{\infty} e^{-ax}\cos(2\pi x\xi)\,dx
=
\frac{a}{a^2+(2\pi\xi)^2},
\]

we get
\[
\hat f(\xi)
=
2\cdot \frac{a}{a^2+(2\pi\xi)^2}.
\]
\pause
\begin{block}{Result}
For $f(x)=e^{-a|x|}$, $a>0$,
\[
\boxed{
\hat f(\xi)=\frac{2a}{a^2+4\pi^2\xi^2}
}
\]

\end{block}
\pause
\begin{itemize}
    \item $\hat f$ is real, positive, even
    \item typical shape of a Lorentzian (Cauchy-type) profile
\end{itemize}

\end{frame}

\subsection{d-dimensional}

\begin{frame}{Fourier Transform in $d$ dimensions}

Let $f:\mathbb{R}^d\to\mathbb{C}$.

\[
\hat f(\xi)
=
\int_{\mathbb{R}^d}
f(x)\,e^{-2\pi i\, x\cdot \xi}\,dx
\]

where
\[
x\cdot \xi = \sum_{j=1}^d x_j\xi_j.
\]
\pause
Inverse transform:

\[
f(x)
=
\int_{\mathbb{R}^d}
\hat f(\xi)\,e^{2\pi i\, x\cdot \xi}\,d\xi.
\]

\end{frame}

\begin{frame}{Main properties in dimension $d$}

\begin{itemize}
    \item Linearity
    \item Translation:
    \[
    f(x-x_0)\ \Longrightarrow\ e^{-2\pi i \xi\cdot x_0}\hat f(\xi)
    \]

    \item Rotation invariance (under orthogonal matrices)
    \[
    f(Qx)\ \Longrightarrow\ \hat f(Q\xi)
    \]

    \item Differentiation:
    \[
    \frac{\partial f}{\partial x_j}
    \ \Longrightarrow\
    2\pi i\xi_j\,\hat f(\xi)
    \]

    \item Laplacian:
    \[
    -\Delta f = - \sum \partial^2_{x_i x_i}f
    \ \Longrightarrow\
    4\pi^2|\xi|^2\hat f(\xi)
    \]
\end{itemize}

\end{frame}

\begin{frame}{Parseval in $\mathbb{R}^d$}

\[
\int_{\mathbb{R}^d} |f(x)|^2\,dx
=
\int_{\mathbb{R}^d} |\hat f(\xi)|^2\,d\xi
\]
\pause
\begin{block}{Meaning}
Energy / variance preserved in any dimension.
\end{block}

Very important for:
\begin{itemize}
    \item PDEs (heat, wave, Schrödinger)
    \item risk decomposition
    \item covariance operators
\end{itemize}

\end{frame}

\begin{frame}{Inversion formula in $\mathbb{R}^d$}

For $f:\mathbb{R}^d\to\mathbb{C}$,

\[
\hat f(\xi)=\int_{\mathbb{R}^d}
f(x)\,e^{-2\pi i x\cdot \xi}\,dx
\]
\pause
\[
f(x)=\int_{\mathbb{R}^d}
\hat f(\xi)\,e^{2\pi i x\cdot \xi}\,d\xi
\]

\begin{block}{Important consequence}
Fourier transform is a unitary operator on $L^2(\mathbb{R}^d)$.
\end{block}

\end{frame}

\begin{frame}{Applications in Financial Risk Engineering}

Multidimensional Fourier transform is used in:

\begin{itemize}
    \item basket option pricing
    \item multi-asset characteristic functions
    \item correlation models
    \item calibration of volatility surfaces
    \item solving multidimensional PDEs
\end{itemize}

\begin{block}{Example}
If $(S_1,\dots,S_d)$ is a vector of log-returns,
its joint characteristic function is a $d$-dimensional Fourier transform.
\end{block}

\end{frame}



\end{document}